\uuid{zuom}
\titre{Notion d'équation différentielle}
\auteur{bodin ; croizat ; sacre}
\organisation{exo7}
\licence{CC BY-NC-SA 4.0}
\datecreate{2022-03-08}
\niveau{L1}
\module{Analyse}
\chapitre{Equation différentielle}
\sousChapitre{Résolution d'équation différentielle du premier ordre}
\difficulte{3}
\theme{Notion d'équation différentielle}
\exo7qcmid{hilisit.equations-differentielles.021}

\contenu{
\begin{question}
\qclassification{225.01}
Soit \(f\) une solution de l'équation différentielle \(y'=2y-x^3\). On sait que la courbe représentative de \(f\) passe par le point \(A(1,2)\). Quelle est la pente de sa tangente au point \(A\) ?
\begin{answers}
	\bad{\(-1\)}
	\bad{\(1\)}
	\bad{\(2\)}
    \good{\(3\)}
\end{answers}
\begin{explanations}
La pente de la tangente au point $A(1,2)$ est le nombre $f'(1)$. Or on sait que \(f(1)=2\) puisque la courbe représentative de \(f\) passe par \(A(1,2)\). De plus, comme \(f\) est solution de l'équation différentielle \(y'=2y-x^3\), on a - en considérant cette égalité pour la fonction $f$ et pour $x=1$ : \(f'(1)=2f(1)-1^3=2\times 2 -1=3\).
\end{explanations}
\end{question}
}
